\section{Introduction}\label{sec:introduction}

An instance of \emph{Subset Sum}, with $n\ge2$, consists of positive integers
$a_1,\ldots,a_n$ and a nonnegative integer $t$.  The question is whether
there is a set $I\subseteq[n]=\{1,\ldots,n\}$ with
$\sum_{i\in I}a_i=t$.  Equal input integers are allowed, and the empty
subset is allowed.  We measure the exponential part of the running time
in the number $n$ of input integers, rather than in their total encoding
length.

Subset Sum is one of the foundational NP-complete problems, appearing
as Knapsack in Karp's list~\cite{Karp1972}.  Its simple formulation has
made it a basic problem in exact exponential-time algorithms and
cryptanalysis.  The meet-in-the-middle algorithm of Horowitz and
Sahni~\cite{HorowitzSahni1974} divides the input into two parts, lists
their subset sums, and searches for complementary sums.  This gives
the familiar $2^{n/2}$ time scale.  Schroeppel and
Shamir~\cite{SchroeppelShamir1981} preserved that time scale while
reducing the space exponent to $1/4$.  More recently, Chen, Jin,
Randolph, and Servedio~\cite{CJRS2023} obtained a randomized worst-case running
time $O(2^{n/2}n^{-\gamma})$ for a constant $\gamma>0.5023$ in word-RAM
and circuit-RAM models with polynomial word length.  A fixed reduction
of the exponential rate below $1/2$ is a different objective from
such polynomial-factor savings.  We
establish that improvement, with ordinary $O(2^{0.49n})$ running time,
in the precise model below.

\subsection{Model and main result}

Let
\[
 b=\max\left\{1,\left\lceil
       \log_2\bigl(1+\max\{t,a_1,\ldots,a_n\}\bigr)
                         \right\rceil\right\}
\]
be the maximum bit length of the input integers and target.
The machine is a word RAM with word length
\begin{equation}\label{eq:word-length}
 w=\left\lceil4\bigl(n+b+\log_2(n+2)\bigr)\right\rceil.
\end{equation}
Each input integer occupies one word in read-only random-access
memory.  Unit-cost operations are word reads and writes, indirect
addressing, comparisons, addition, subtraction, multiplication,
integer quotient and remainder with a nonzero divisor, bitwise Boolean
operations, and shifts.  Each operation has a constant number of
$w$-bit inputs and outputs.  Arithmetic exceeding these limits is
implemented by multiple word operations and charged accordingly.
A fresh independent uniform $w$-bit random word costs one operation.
Randomness is one-way: an earlier random word can be reused only if it
has been retained in writable memory.  There is no random oracle,
uncharged external storage, or separate workspace bound.

The algorithm must be a single finite program for all $n$ and $b$,
without advice or externally supplied tables depending on $n$ or the
input.  All preprocessing, table construction, input accesses, and
random sampling are included in its running time.

\begin{theorem}\label{thm:main}
There is a uniform randomized classical algorithm $\mathcal A$ for
Subset Sum with the following properties in the word-RAM model
\eqref{eq:word-length}.
For every input with $n\ge2$, the algorithm halts on every execution
and outputs the correct answer with probability at least $2/3$.
For every fixed positive integer $c$, there are constants $C_c>0$
and $N_c\ge2$ such that every input with $n\ge N_c$ and $b\le n^c$
is processed, on every outcome of the random choices, in at most
\[
                         C_c\,2^{0.49n}
\]
word operations.  The program is not given $c$ and does not depend
on $c$.
\end{theorem}

The probability in the theorem is over the algorithm's internal
random words, for each fixed input.  In particular, no distributional
assumption is made on the input integers.  The running-time bound
applies even to executions on which the answer is incorrect.
Our intermediate bound has the form
$2^{0.489995n}\operatorname{poly}(n+b+2)$; the strict exponent slack
absorbs the polynomial factor separately for every fixed $c$.

\subsection{Related approaches}

Two other directions explain both the progress on Subset Sum and the
difficulty of a worst-case improvement in $n$.  Pseudopolynomial
algorithms measure their complexity in the numerical target $t$.
Bellman's dynamic-programming approach~\cite{Bellman1957} yields
$O(nt)$ time.  Koiliaris and Xu~\cite{KoiliarisXu2017} improved this
bound deterministically, and Bringmann~\cite{Bringmann2017} obtained
a randomized $\widetilde O(n+t)$ algorithm.  Chan~\cite{Chan2026}
subsequently obtained a deterministic near-linear pseudopolynomial
algorithm.  These bounds are particularly effective for small targets;
polynomial \emph{bit length}, however, permits much larger targets.
Here and in historical comparisons, $\widetilde O$ suppresses
polylogarithmic factors.

Representation methods provide a different kind of saving.  The
method introduced by Howgrave-Graham and Joux~\cite{HGJ2010}, and
developed by Becker, Coron, and Joux~\cite{BCJ2011}, represents one
solution by many pairs of partial objects and then searches a filtered
fraction of those representations.  Their improvements concern
random instances, not a guarantee for every input.  Austrin, Kaski,
Koivisto, and Nederlof~\cite{AKKN2016} developed algorithms exploiting
input blocks with many distinct sums through representations, or
blocks with few distinct sums through faster enumeration of those
sums.  Nederlof and
W\k{e}grzycki~\cite{NederlofWegrzycki2021} used related ideas to improve
the space bound below the Schroeppel--Shamir exponent.

Related work concerns \emph{Subset Balancing}: given integers $v_i$
and a fixed finite coefficient set $C\subset\mathbb Z$, find a nonzero
$x\in C^n$ satisfying $\sum_i v_ix_i=0$.
Randolph and W\k{e}grzycki~\cite{RandolphWegrzycki2026} give worst-case
exponent savings for several coefficient sets, though not for ordinary
Subset Sum.  Their compatibility certificates recover valid
representation pairs without inspecting all invalid pairs.
Our alias extraction uses the representation and distinct-sum
perspective for a different task: listing modular false positives so
that their contributions can be removed from an estimated checksum.

\subsection{Proof overview}

The proof separates two complementary tasks:
estimating a modular target bin and identifying the false positives
created by the modulus.  First, random tie-breaking weights transform
the original problem into polynomially many targets for new positive
integer weights $c_1,\ldots,c_n$.  The isolation lemma of Mulmuley,
Vazirani, and Vazirani~\cite{MVV1987} ensures that, with high constant
probability, some transformed target has exactly one solution whenever
the original instance is positive.  A preliminary meet-in-the-middle
test handles inputs whose partial subset sums have sufficiently small
support.  Both reductions are developed in Section~\ref{sec:setup}.

For a prime $P$, write $e_P(x)=\exp(2\pi\mathrm i x/P)$.
With independent uniform residues
$\phi_1,\ldots,\phi_n$ modulo $P$, assign a subset $I$ the phase
\[
 \chi(I)=e_P\left(\sum_{i\in I}\phi_i\right).
\]
For each transformed target $t'$, a subset is an \emph{alias} if its
sum is $t'$ modulo $P$ but not over the integers.  The identity
\begin{equation}\label{eq:overview-subtraction}
 \sum_{\sum_{i\in I}c_i=t'}\chi(I)
  =\sum_{\sum_{i\in I}c_i\equiv t'\pmod P}\chi(I)
     -\sum_{\substack{I\text{ an alias}\\\text{for }t'}}\chi(I)
\end{equation}
is the organizing principle.  On a negative instance the left side is
zero for every transformed target; an isolated positive target leaves
a single phase of modulus one.  The two sums on the right can be much
larger, so their difference must be recovered to small \emph{absolute}
error.  A relative approximation to the modular sum would not suffice.
We choose a random prime modulus larger than the desired running time,
so neither scanning all residues nor evaluating every Fourier frequency
is affordable.  We estimate the modular sum and subtract the phases of
a complete list of aliases.

\paragraph{A checksum without a full Fourier transform.}
Writing $V_{t'}$ for the modular sum in
\eqref{eq:overview-subtraction}, Fourier inversion gives
\[
 V_{t'}=\frac1P\sum_{k=0}^{P-1}J_{t'}(k),
 \qquad
 J_{t'}(k)=e_P(-kt')\prod_{i=1}^n
                         \bigl(1+e_P(\phi_i+kc_i)\bigr).
\]
For each fixed $k$, independent uniform phases give
$\mathbb E_\phi|J_{t'}(k)|^2=2^n$.  Thus the elementary second-moment
bound for uniform frequency sampling does not give the required
running time.
For a chosen digit base $L$, writing $k=hL+l$ expresses the phase
vector $(\phi_i+kc_i)_i$ as the difference between
$(\phi_i+hLc_i)_i$ and $(-lc_i)_i$, all modulo $P$.
These high- and low-digit vectors form two lists.  The filters must
keep the stored lists short, cover every frequency, and control the
second moment of the resulting estimate.  A table specifies
random coordinate shifts, with some coordinates ignored; a sum of
logarithmic density scores tests each vector against the table.
A \emph{match} consists of a table and a digit pair whose two vectors
both pass.  This two-sided organization is related to asymmetric
locality-sensitive filtering~\cite{Christiani2017}.

To analyze coverage of a fixed digit pair, we temporarily change the
random-shift law at each coordinate.  Three bounded functions record
the two resulting mean scores and the likelihood-ratio cost of that change.
A \emph{pattern} consists of coarse averages of these functions on
coordinate blocks; the tests are chosen separately for each pattern.
Grouping coordinates into blocks makes the number of patterns and the
cost of generating block tables subexponential.  Comparing this auxiliary
law with the actual table law shows that, with high probability, every
frequency appears in a match for its own pattern.
Two fixed inequalities for circle densities control the second
moment.  Section~\ref{sec:filtering} states these inequalities, and
Appendix~\ref{app:numerical} proves them with rational enclosures and
analytic quadrature bounds.

Patterns with a sufficiently small high-probability bound on their
total match count are summed directly.  For the other patterns, the
algorithm samples matches and
divides by the number of matches representing the sampled frequency.
This is the inverse-coverage correction familiar from union-size
estimation~\cite{KLM1989}, applied here to a complex weighted sum.
Theorem~\ref{thm:checksum} proves the resulting absolute-error estimate
for arbitrary input residues, independently of the distinct-sum
condition used in alias recovery.  Random phases make the coordinates
independent at each fixed frequency, even though the original input
is arbitrary.  The second-moment calculation uses the original phase
law, not the law conditioned on successful filtering.

\paragraph{Extracting only aliases.}
The second procedure, in Section~\ref{sec:aliases}, gives each subset
many representations by two records.  After randomly complementing
some subset-indicator variables, choose a shared coordinate block and
partition its complement into a left and a right side.  Each record
consists of a subset of its side together with a \emph{mask}, a subset
of the shared block.  Disjoint masks give an actual subset;
overlapping masks count their common coordinates twice in the sum of
the record weights.  A small auxiliary modulus reduces the record
lists.  The analysis proves a dichotomy: either the preliminary
distinct-sum test succeeds with high probability, or enough
representations survive to recover every alias with high probability.
The enumeration omits every pair satisfying the corresponding exact
integer equation, including pairs with overlapping masks.  This
omission is crucial: each remaining pair has a nonzero integer
discrepancy, so a random-prime divisor bound controls the number of
pairs that pass the main modulus.  Disjoint masks are checked only
after this bounded enumeration.

Once the reconstructed left side of \eqref{eq:overview-subtraction}
has error less than $1/4$, comparison with the threshold $1/2$
separates an empty bin from an isolated one.  Thus isolation removes any need for a separate
anti-concentration theorem for sums of phases.
Section~\ref{sec:implementation} assembles the procedures and proves
Theorem~\ref{thm:main}, including finite-precision arithmetic, capped
random sampling, and a running-time bound on every execution.
